\documentclass[11pt,a4paper]{article}

% ------- Packages ---------------------------------------------------------
\usepackage[T1]{fontenc}
\usepackage[utf8]{inputenc}
\usepackage[margin=1in]{geometry}
\usepackage{amsmath,amssymb}
\usepackage{booktabs}
\usepackage{graphicx}
\usepackage{xcolor}
\usepackage{listings}
\usepackage{hyperref}
\hypersetup{colorlinks=true,linkcolor=blue!60!black,urlcolor=blue!60!black,citecolor=blue!60!black}

% ------- Code listing style ----------------------------------------------
\definecolor{codebg}{RGB}{248,248,248}
\definecolor{keyword}{RGB}{0,102,153}
\definecolor{string}{RGB}{170,55,55}
\definecolor{comment}{RGB}{100,100,100}

\lstdefinestyle{py}{
    language=Python,
    backgroundcolor=\color{codebg},
    basicstyle=\ttfamily\small,
    keywordstyle=\color{keyword}\bfseries,
    stringstyle=\color{string},
    commentstyle=\color{comment}\itshape,
    showstringspaces=false,
    breaklines=true,
    frame=single,
    framesep=4pt,
    xleftmargin=0pt,
    numbers=left,
    numberstyle=\tiny\color{gray},
    numbersep=8pt,
    tabsize=4
}
\lstset{style=py}

\lstdefinestyle{sh}{
    language=bash,
    backgroundcolor=\color{codebg},
    basicstyle=\ttfamily\small,
    keywordstyle=\color{keyword},
    commentstyle=\color{comment}\itshape,
    showstringspaces=false,
    breaklines=true,
    frame=single,
    framesep=4pt,
    xleftmargin=0pt
}

% ------- Title ------------------------------------------------------------
\title{\textbf{ViT-JAX Distributed Scaling}\\
       \large Study notes on data-parallel Vision Transformer training on TPU}
\author{Project notes}
\date{}

\begin{document}
\maketitle

\begin{abstract}
\noindent This document is a study-notes writeup of the \texttt{vit-jax-distributed-scaling}
project: we train a Vision Transformer (ViT-Small) on CIFAR-100 using JAX/Flax, parallelised
with \texttt{jax.pmap}, and use the setup to probe two systems questions:
(i) \emph{how does throughput scale with device count?} and
(ii) \emph{how badly does a single slow worker hurt synchronous training?}
Every section pairs the conceptual intuition with the minimal piece of code that implements
it, so the document doubles as a guided reading of the repository.
\end{abstract}

\tableofcontents
\vspace{1em}
\hrule
\vspace{1em}

% =========================================================================
\section{What we built, and why}
% =========================================================================

A Vision Transformer trains by (a) splitting an image into patches,
(b) treating the patches as a token sequence, and (c) feeding that sequence through a stack of
transformer blocks. Training one is cheap; training one \emph{at scale, across accelerators}
exposes two systems phenomena that matter far beyond ViT:

\begin{itemize}
\item The \textbf{all-reduce barrier}: every device must reach the gradient-average step before any
      can proceed. This is the single most important synchronisation primitive in data-parallel
      training, and it is what makes distributed training brittle to slow workers.
\item \textbf{Scaling efficiency}: doubling the devices never exactly doubles throughput because
      gradients have to travel over an interconnect. Measuring how much you lose is how you tell
      whether your distributed strategy is working.
\end{itemize}

We use ViT-Small on CIFAR-100 because the model is small enough ($\sim$10M parameters) that
per-step work is short, so the all-reduce and straggler effects dominate the signal we want to
measure. On a bigger model the compute would hide everything.

\subsection{Repo layout}

\begin{lstlisting}[style=sh]
vit_jax_distributed/
  models/vit.py             # ViT-Small (Flax Linen)
  data/cifar100.py          # TFDS -> numpy in-memory pipeline
  distributed/parallel.py   # pmap, pmean, train state
  train/trainer.py          # training loop
  experiments/
    scaling.py              # throughput vs device count
    straggler.py            # step time with a slow worker
  utils/                    # config, logging, timing
main.py                     # CLI entry point
scripts/setup.sh            # installs jax[tpu|cuda|cpu] per host
scripts/run_gcp_tpu.sh      # end-to-end TPU run
\end{lstlisting}

\subsection{Smoke-test result (the numbers we are going to explain)}

On an 8-chip TPU v6e spot VM, a two-epoch run produced:

\begin{center}
\begin{tabular}{lr}
\toprule
Model parameters               & 9{,}554{,}404 \\
Global batch size              & 512 \\
Per-device batch size          & 64  \\
Steady-state step time         & 32.8 ms (p50) \\
First step (XLA compile)       & $\sim$65 s \\
Post-compile throughput        & $\sim$15{,}600 images/sec \\
Test accuracy after 2 epochs   & 22.0\% \\
\bottomrule
\end{tabular}
\end{center}

The 65-second first step is not a bug; it is XLA tracing and compiling the pmap'd training
step into a TPU executable. Every subsequent step re-uses that compiled graph, which is why
they are $\sim$2000$\times$ faster.

% =========================================================================
\section{The model: ViT-Small}
% =========================================================================

\subsection{Intuition: images as sequences}

Transformers are sequence models. To feed them an image we turn the image into a sequence:

\begin{enumerate}
\item Tile the $32 \times 32$ image into non-overlapping $4 \times 4$ patches, yielding 64 patches.
\item Flatten each patch and project it through a learnable linear layer to get a 384-dim token.
\item Prepend a learnable \texttt{[CLS]} token (so we have 65 tokens total).
\item Add learnable positional embeddings (the model cannot recover spatial position otherwise).
\end{enumerate}

The output of this preprocessing is a tensor of shape $(B, 65, 384)$: a batch of 65-token
sequences, each token a 384-dim embedding. From here on it is pure Transformer encoder.

\subsection{Pre-norm blocks}

Each transformer block is:
$$
x' = x + \text{Attn}(\text{LN}(x)), \qquad x'' = x' + \text{MLP}(\text{LN}(x'))
$$

Applying LayerNorm \emph{before} attention and MLP (``pre-norm'') is the standard choice in modern
ViT and GPT-style models: it stabilises training at depth because the residual stream stays at
unit scale regardless of how many layers you stack.

\subsection{Classification head}

After $N$ transformer blocks and a final LayerNorm, we read out the \texttt{[CLS]} token only
(\texttt{x[:, 0]}) and feed it through a linear layer to 100 logits. The intuition: the
self-attention has, over $N$ rounds, aggregated information from every patch \emph{into} the CLS
token, so it is a learned summary of the whole image.

\subsection{Code: patch embedding}

\begin{lstlisting}[style=py]
class PatchEmbedding(nn.Module):
    hidden_dim: int = 384
    patch_size: int = 4
    image_size: int = 32

    @nn.compact
    def __call__(self, x):
        B = x.shape[0]
        N = (self.image_size // self.patch_size) ** 2    # 64 patches
        # (B, H, W, C) -> (B, N, patch_size*patch_size*C) -> (B, N, hidden_dim)
        x = nn.Dense(self.hidden_dim, name="patch_projection")(
            x.reshape(B, self.image_size // self.patch_size, self.patch_size,
                         self.image_size // self.patch_size, self.patch_size, -1)
             .transpose(0, 1, 3, 2, 4, 5)
             .reshape(B, N, -1))
        cls = self.param("cls_token",    nn.initializers.normal(0.02), (1, 1, self.hidden_dim))
        pos = self.param("pos_embedding", nn.initializers.normal(0.02), (1, N + 1, self.hidden_dim))
        x   = jnp.concatenate([jnp.broadcast_to(cls, (B, 1, self.hidden_dim)), x], axis=1)
        return x + pos
\end{lstlisting}

Full architecture (\texttt{vit\_jax\_distributed/models/vit.py}): 8 transformer blocks, hidden
384, 6 heads, MLP dim 768, $\sim$9.5M parameters. Standard ViT-Small-for-CIFAR.

% =========================================================================
\section{Data pipeline}
% =========================================================================

CIFAR-100 is 50{,}000 training images, $32 \times 32 \times 3$ each. Standardised and stored as
float32 that is $50000 \cdot 32 \cdot 32 \cdot 3 \cdot 4 \approx 614$ MB. The full dataset fits
comfortably in RAM, so we skip \texttt{tf.data} entirely and use a plain numpy iterator:

\begin{lstlisting}[style=py]
def _numpy_data_iterator(images, labels, batch_size, shuffle, augment, loop, seed):
    n     = images.shape[0]
    usable = (n // batch_size) * batch_size          # drop the ragged tail
    rng    = np.random.RandomState(seed)
    while True:
        idx = np.arange(n)
        if shuffle: rng.shuffle(idx)
        idx = idx[:usable]
        for s in range(0, usable, batch_size):
            batch_idx  = idx[s : s + batch_size]
            batch_imgs = images[batch_idx].copy()
            if augment:
                batch_imgs = _augment_batch(batch_imgs, rng)
            yield {"image": batch_imgs, "label": labels[batch_idx]}
        if not loop: break
\end{lstlisting}

\textbf{Augmentation.} Random horizontal flip and random crop with 4-pixel reflection padding.
Both are vectorised with advanced indexing over the whole batch; no per-image Python loop.

\subsection{Sharding for pmap}

Every distributed step needs each device's private slice of the batch. We reshape
$(B, H, W, C) \to (\text{num\_devices}, B/\text{num\_devices}, H, W, C)$ and let pmap dispatch
slice $i$ to device $i$:

\begin{lstlisting}[style=py]
def shard_batch(batch, num_devices):
    def reshape(x):
        per_device = x.shape[0] // num_devices
        return x.reshape((num_devices, per_device) + x.shape[1:])
    return {k: reshape(v) for k, v in batch.items()}
\end{lstlisting}

A global batch of 512 across 8 devices means each device sees 64 images per step.

% =========================================================================
\section{Distributed training: the conceptual heart}
% =========================================================================

\subsection{SPMD and pmap}

JAX's \texttt{jax.pmap} is a SPMD (``Single Program, Multiple Data'') primitive. You write one
function; pmap arranges for it to run on every device simultaneously, each device receiving its
own slice of the input.

\begin{lstlisting}[style=py]
@partial(jax.pmap, axis_name="batch", devices=devices)
def train_step(state, batch, rng):
    ...
\end{lstlisting}

Two things worth internalising:
\begin{itemize}
\item Inputs to a pmap'd function must have a leading dimension of size \texttt{len(devices)};
      pmap strips this dimension and gives device $i$ the $i$-th slice.
\item Outputs \emph{keep} a leading device dimension. To read a scalar loss on the host you
      usually index \texttt{loss[0]} --- because loss was \texttt{pmean}'d across devices, every
      device holds the same value.
\end{itemize}

\subsection{All-reduce via pmean}

The magic happens inside the step: after each device has computed its local gradients from its
slice of the batch, we \emph{average} those gradients across every device. That averaging is the
\textbf{all-reduce} operation, and it is what makes the whole thing mathematically equivalent to
training on one device with the full global batch.

\begin{lstlisting}[style=py]
grads = jax.lax.pmean(grads, axis_name="batch")   # <-- the all-reduce
loss  = jax.lax.pmean(loss,  axis_name="batch")
acc   = jax.lax.pmean(acc,   axis_name="batch")
new_state = state.apply_gradients(grads=grads)
\end{lstlisting}

The name \texttt{"batch"} is arbitrary; it names the pmap axis. You can have nested pmaps with
different axis names (``batch'' for data parallel, ``model'' for tensor parallel) --- not used in
this project, but the pattern generalises.

\subsection{Why this matters: the barrier}

All-reduce is a \textbf{collective} operation: it cannot complete on any single device until
every device has contributed its gradient. Every training step therefore has an implicit barrier
where slow workers hold up fast ones. Schematically:

\begin{lstlisting}[style=sh]
Device 0: forward -> backward -> [pmean] -> optimizer step
Device 1: forward -> backward -> [pmean] -> optimizer step
Device 2: forward -> backward -> [pmean] -> optimizer step
Device 3: forward -> backward -> [pmean] -> optimizer step
                                    ^
                           all-reduce barrier:
                           nobody moves until everybody arrives
\end{lstlisting}

This is the phenomenon the \textbf{straggler experiment} isolates.

\subsection{Train state and replication}

Flax's \texttt{TrainState} bundles model apply function + parameters + optimizer state. We
create one master copy, then replicate it across every device with \texttt{flax.jax\_utils.replicate},
which prepends a device axis of size \texttt{len(devices)}:

\begin{lstlisting}[style=py]
def create_train_state(rng, model, lr, image_size, weight_decay, warmup_steps, total_steps):
    dummy  = jnp.ones([1, image_size, image_size, 3])
    params = model.init(rng, dummy, train=False)["params"]

    effective_warmup = min(warmup_steps, max(1, total_steps // 2))
    schedule = optax.warmup_cosine_decay_schedule(
        init_value=0.0, peak_value=lr,
        warmup_steps=effective_warmup, decay_steps=total_steps, end_value=0.0)
    tx = optax.adamw(learning_rate=schedule, weight_decay=weight_decay)
    return TrainState.create(apply_fn=model.apply, params=params, tx=tx)

state = replicate_state(state, devices=devices)     # <- broadcast to all devices
\end{lstlisting}

The clamp on \texttt{effective\_warmup} exists because \texttt{optax.warmup\_cosine\_decay\_schedule}
internally computes \texttt{decay\_steps - warmup\_steps} and errors out on a negative result ---
which is exactly what happens when you run a two-epoch smoke test with the default 500-step
warmup. Clamping to $\text{total\_steps}/2$ keeps the long-run behaviour unchanged while making
short runs legal.

% =========================================================================
\section{Putting it together: the training loop}
% =========================================================================

\begin{lstlisting}[style=py]
state        = replicate_state(state, devices=devices)
p_train_step = create_train_step(devices=devices)
p_eval_step  = create_eval_step(devices=devices)
train_iter, _ = get_datasets(config, num_devices=num_devices, seed=seed)

for epoch in range(config.epochs):
    for step in range(steps_per_epoch):
        batch        = next(train_iter)
        batch        = shard_batch(batch, num_devices)    # (N, B/N, ...)
        rng, step_rng = jax.random.split(rng)
        step_rngs    = jax.random.split(step_rng, num_devices)

        state, metrics = p_train_step(state, batch, step_rngs)

        step_loss = float(metrics["loss"][0])    # implicit host barrier
        step_acc  = float(metrics["accuracy"][0])

    # end of epoch: evaluate
    test_metrics = _run_evaluation(state, p_eval_step, config, num_devices)
\end{lstlisting}

Four things worth noticing:

\begin{enumerate}
\item \textbf{One RNG per device.} \texttt{jax.random.split(step\_rng, num\_devices)} gives each
      device its own dropout stream. Without this, every device would apply identical dropout,
      reducing regularisation strength.
\item \textbf{Asynchronous dispatch.} \texttt{p\_train\_step} returns immediately with lazy
      arrays; the compute runs on TPU in the background. The \texttt{float(metrics["loss"][0])}
      call is what blocks the host until the step is actually done --- this is our implicit
      barrier for timing.
\item \textbf{Metrics are pmean'd inside the pmap}, so \texttt{metrics["loss"]} has shape
      \texttt{(num\_devices,)} and every entry is equal. Indexing \texttt{[0]} just picks any one.
\item \textbf{Per-step batch sharding.} Sharding is a host-side numpy reshape, not a TPU op, so
      it is practically free compared to the $\sim$33 ms step itself.
\end{enumerate}

% =========================================================================
\section{Experiments}
% =========================================================================

\subsection{Scaling experiment}

\subsubsection*{What we are measuring}

For $k \in \{1, 2, 4, 8\}$ we measure:
\begin{itemize}
\item \textbf{Step time} $T_k$: median wall-clock time to execute one forward-backward-update
      cycle on $k$ devices with per-device batch size held constant, so the global batch grows
      with $k$.
\item \textbf{Throughput} $\Theta_k = B_k / T_k$ images per second, where $B_k = k \cdot B_1$.
\item \textbf{Scaling efficiency} $\eta_k = \Theta_k / (k \cdot \Theta_1)$. Perfect scaling is
      $\eta = 1$.
\end{itemize}

\subsubsection*{What to expect, and why}

Gradient averaging requires sending parameter-sized tensors between devices. Total all-reduce
bandwidth is fixed, so as $k$ grows the constant communication cost takes a larger \emph{relative}
bite out of step time. Concretely, if per-step compute is $C$ and all-reduce is $R$,
$$
T_k \approx C + R, \qquad \eta_k \approx \frac{C}{C + R}
$$
(ignoring topology, overlap, and the fact that $R$ itself grows with $k$). For our tiny model
$C \sim 30$ ms and $R$ is small on a TPU v6e mesh, so we expect $\eta \gtrsim 0.85$ at $k=8$.

\subsubsection*{Measurement hygiene}

\begin{itemize}
\item \textbf{Warmup.} The first step compiles via XLA ($\sim$65 s on v6e-8). We run 10 warmup
      steps before we start timing.
\item \textbf{Explicit block.} We call \texttt{jax.tree.map(lambda x: x.block\_until\_ready(), state)}
      after each timed step, so the tick captures actual device work and not just dispatch latency.
\item \textbf{Device subset.} To benchmark a strict subset of devices, we forward
      \texttt{devices=all\_devices[:k]} into \emph{both} \texttt{jax.pmap} \emph{and}
      \texttt{flax.jax\_utils.replicate}. Forgetting either leaves pmap defaulting to all local
      devices, and the sharded batch's leading dimension then mismatches.
\end{itemize}

\subsection{Straggler experiment}

\subsubsection*{The question}

A synchronous data-parallel system only advances as fast as its slowest worker. How much does a
single slow device hurt total throughput?

\subsubsection*{Injecting a real straggler}

Naive \texttt{time.sleep(x)} is useless here: it is a Python call, and Python does not run
\emph{inside} a pmap'd, JIT-compiled XLA program. Even if it did, an intelligent scheduler
might just reorder around it. We need a delay XLA cannot elide, which means real compute with a
real data dependency on itself:

\begin{lstlisting}[style=py]
device_idx = jax.lax.axis_index("batch")          # 0..num_devices-1

def _straggler(acc):
    def body(i, a): return a @ jnp.eye(64) + 0.0 * a
    return jax.lax.fori_loop(0, delay_iterations, body, acc)

def _no_delay(acc): return acc

dummy = jnp.ones((64, 64))
_ = jax.lax.cond(device_idx == 0, _straggler, _no_delay, dummy)
\end{lstlisting}

Each loop iteration multiplies a $64 \times 64$ matrix by the identity and adds
$0 \cdot a$: the data dependency on the previous iteration's \texttt{acc} stops the compiler
from constant-folding or hoisting the loop, and the \texttt{lax.cond} on
\texttt{axis\_index("batch") == 0} confines the work to one device.

\subsubsection*{What to expect}

For baseline step time $T$ and injected extra time $\Delta$ on one device,
$$
T_{\text{straggler}} \approx T + \Delta, \qquad
\text{slowdown} = \frac{T_{\text{straggler}}}{T} \approx 1 + \frac{\Delta}{T}.
$$
In other words: \emph{every} device pays the straggler tax even though only one device is slow.
A single bad worker in a 1000-GPU cluster makes the other 999 wait.

The experiment sweeps \texttt{delay\_iterations} $\in \{100, 500, 1000, 5000, 10000\}$ and
reports the ratio of straggler step time to baseline step time. The curve on log-$x$ axis should
be roughly linear once $\Delta$ is large enough to dominate $T$.

\subsubsection*{Why this matters in production}

Real clusters have real stragglers: thermal throttling, bad NICs, GC pauses, a noisy neighbour
on a shared node. The straggler experiment is a miniature of why techniques like
asynchronous SGD, elastic averaging, hedged requests, and backup workers exist.

% =========================================================================
\section{TPU-specific notes}
% =========================================================================

\subsection{Why JAX installs differently on TPU}

TPU does not speak CUDA; it speaks over \texttt{libtpu}, a Google-provided shared library that
plugs into JAX via a plugin. You \textbf{cannot} \texttt{pip install jax[cuda12]} on a TPU VM ---
JAX will import successfully but fall back to CPU, silently, and your training will be
$\sim$50$\times$ slower with no error. The correct install command is:

\begin{lstlisting}[style=sh]
pip install -U "jax[tpu]>=0.4.20" \
    -f https://storage.googleapis.com/jax-releases/libtpu_releases.html
\end{lstlisting}

\texttt{scripts/setup.sh} detects TPU hosts by checking \texttt{/dev/accel*} and the
\texttt{TPU\_NAME} env var --- deliberately \emph{before} JAX is installed, so it can choose the
right wheel on a fresh VM.

\subsection{bf16 on the MXU}

TPU v3 and v6e compute matmuls natively in bf16 (brain-float: 1 sign bit, 8 exponent bits,
7 mantissa bits). Enabling bf16 for matmuls while keeping fp32 parameters is a one-liner:

\begin{lstlisting}[style=py]
jax.config.update("jax_default_matmul_precision", "bfloat16")
\end{lstlisting}

This is what our \texttt{--precision bf16} flag does. It does \emph{not} change parameter dtype
or optimizer state; it only tells XLA that every matmul (dense layers, attention QKV, einsum
inside convolutions) can run the MXU in bf16 mode. On v6e this is typically a $1.5$--$2\times$
throughput win for essentially free --- the narrower matmul precision is absorbed by the
fp32 accumulator and the larger numbers that dominate training.

\subsection{pmap vs pjit}

This project uses \texttt{pmap} exclusively, because we only ever do \emph{pure data
parallelism}: the same model on every device. \texttt{pjit} (and its successor
\texttt{shard\_map}) is the right tool when you need \emph{tensor parallelism} (a single layer
split across devices) or \emph{pipeline parallelism}. Our ViT-Small fits in one TPU chip's
HBM, so model parallelism would be gratuitous.

\subsection{Spot/preemption}

v6e-8 spot VMs cost roughly one third of on-demand (\$1.26/hr vs \$4.21/hr as of this writing)
but can be \emph{preempted} at any time: GCP reclaims the hardware, the VM transitions to
\texttt{PREEMPTED}, and your process is gone. This hit us during initial setup, once. Mitigation
in practice: either use on-demand for long runs, or use \texttt{gcloud alpha compute tpus
queued-resources} to auto-retry until capacity returns.

% =========================================================================
\section{Running it}
% =========================================================================

Local laptop (CPU) or single GPU:
\begin{lstlisting}[style=sh]
bash scripts/setup.sh                    # detects hardware, installs jax build
python main.py --experiment train --batch_size 128 --epochs 2
\end{lstlisting}

GCP TPU v6e-8:
\begin{lstlisting}[style=sh]
gcloud compute tpus tpu-vm create vit-training \
    --zone=us-central1-b --accelerator-type=v6e-8 \
    --version=v2-alpha-tpuv6e --spot

gcloud compute tpus tpu-vm scp --recurse --zone=us-central1-b \
    main.py requirements.txt scripts vit_jax_distributed \
    vit-training:~/vit-jax-distributed-scaling/

gcloud compute tpus tpu-vm ssh vit-training --zone=us-central1-b \
    --command="cd ~/vit-jax-distributed-scaling && bash scripts/setup.sh"

gcloud compute tpus tpu-vm ssh vit-training --zone=us-central1-b \
    --command="cd ~/vit-jax-distributed-scaling && bash scripts/run_gcp_tpu.sh"

# Remember to delete afterwards --- spot VM keeps billing until you do.
gcloud compute tpus tpu-vm delete vit-training --zone=us-central1-b
\end{lstlisting}

% =========================================================================
\section{What to look at in the outputs}
% =========================================================================

Every run writes three files to \texttt{--output\_dir}:

\begin{itemize}
\item \texttt{step\_metrics.csv} --- one row per logged step: loss, accuracy, step time.
\item \texttt{epoch\_metrics.csv} --- one row per epoch: aggregated train + test metrics.
\item \texttt{training\_log.json} --- config, metadata, and all step/epoch logs.
\end{itemize}

Each experiment adds plots:

\begin{itemize}
\item Scaling: \texttt{throughput\_vs\_devices.png}, \texttt{step\_time\_vs\_devices.png},
      \texttt{scaling\_efficiency\_vs\_devices.png}.
\item Straggler: \texttt{step\_time\_comparison.png}, \texttt{slowdown\_vs\_delay.png}.
\end{itemize}

The scaling-efficiency plot is the most interpretable single number: a flat line near $1.0$
means your distributed setup is not losing throughput to communication; a downward slope means
it is.

% =========================================================================
\section{Further reading and things to try}
% =========================================================================

\begin{itemize}
\item \textbf{Swap CIFAR for something bigger.} Imagenette or Tiny-ImageNet will push per-step
      compute higher, making the all-reduce overhead relatively smaller and the scaling
      efficiency closer to $1.0$. A useful contrast to run.
\item \textbf{Try \texttt{pjit}/\texttt{shard\_map}.} Port the same model to \texttt{jit} + a
      SPMD mesh; you will see near-identical throughput for pure data parallelism but cleaner
      code.
\item \textbf{Compute-communication overlap.} Look at what happens if you replace
      \texttt{pmean(grads)} with a manual \texttt{psum} and start the all-reduce asynchronously
      during the optimizer step.
\item \textbf{Real straggler sources.} Replace the synthetic \texttt{fori\_loop} with a
      \texttt{jax.lax.while\_loop} keyed on a random per-device delay, to see the effect of
      \emph{variance} in per-device times rather than a constant straggler.
\end{itemize}

\end{document}
